\section{Introduction}
\label{sec:introduction}

In the history of economic thought, progress rarely unfolds as a linear, cumulative ascent toward universal truth. As \citet{Roncaglia2006} observes, theoretical advance throughout the history of economic thought is governed by \textit{competing views}: an enduring contest among rival paradigms rooted in fundamentally distinct conceptualizations, analytical categories, and contested world visions. For over half a century, the history of Chile's Unidad Popular (UP, 1970--1973) under Salvador Allende has been interpreted through an uncritical cumulative lens hegemonized by mainstream economics \citep{LarrainMeller1991,Edwards2024}. Mainstream accounts under neoclassical market-clearing assumptions and monetarist approaches to the balance of payments and inflation that preclude distributive conflict, have treated the Chilean breakdown as a cautionary tale of populist fiscal mismanagement. Yet beneath this orthodox consensus lies an unexamined failure of conceptualization: treating peripheral capitalism as a closed monetary system, blind to the structural subordination of dependent economies within an international currency hierarchy marked by external balance-of-payments constraints and financial subordination \citep{Jessop2019,Chamorro2025}.

This chapter investigates the causal ordering between monetary creation and inflation, and how central-bank foreign-exchange solvency mediates the transmission of domestic liquidity during a peripheral socialist transformation subject to an external constraint. If domestic money creation mechanically drove inflation, why did Central Bank base money expand by over 100\% in 1971 while monthly inflation remained quiescent at 1.5--2.0\% and annual inflation fell from 35\% to 22\%? The answer lies in balance-sheet solvency and real capacity: in 1971, the Central Bank possessed nearly \$400~million in inherited gross reserves and industry mobilized significant idle capacity, with structural capacity utilization surging to 98.4\%, absorbing liquidity through transactions demand. Significant inflation pass-through did not take place in 1971; it emerged only in mid-1972 once international reserves were depleted and a targeted external credit blockade severed revolving commercial trade lines. Domestic money growth operated as an endogenous accommodation valve through which the monetary authority mediated an acute balance-of-payments squeeze amidst intense distributive conflict. Contesting both the monetarist thesis of \citet{Edwards2024} and the Austrian interpretation of \citet{EspinosaVejar2025}, I formalize the central-bank balance sheet through the Solvency Ratio---the ratio of convertible foreign reserve assets to domestic high-powered liabilities ($S_t \equiv F_t / H_t$)---and evaluate non-linear transmission using the monthly Solvency Growth Gap ($\Delta s_{t-1} \equiv g^F_{t-1} - g^H_{t-1}$). Threshold econometric estimates identify a structural transition into conditions of central bank insolvency at $\Delta s_{t-1} \le -4.615\%$. Within this insolvency regime, generalized impulse response functions establish that forward monetary transmission ($g^H \to \pi_m$) remains statistically significant for 9 consecutive months, delivering a cumulative 12-month price increase of 3.55 percentage points. In contrast, reverse accommodation ($\pi_m \to g^H$) remains statistically significant for 10 consecutive months, yielding a cumulative 12-month base-money expansion of 6.54 percentage points. Reverse accommodation significantly exceeds forward transmission in both persistence and cumulative magnitude, with point-wise multiplier differences significant across 9 consecutive months ($h=2,\dots,10$). Monetary transmission is thus state-dependent on external reserve solvency.

Allende's thousand days unfolded within the tectonic shifts of the global structural crisis of Fordist capitalism during the early 1970s. As \citet{Zoeller2019} demonstrates, President Richard Nixon's unilateral suspension of dollar-gold convertibility on August 15, 1971, was not an inevitable systemic accident, but a contingent political-economic maneuver undertaken to preserve U.S. domestic policy autonomy amidst growing balance-of-payments pressures. While dismantling the Bretton Woods par-value system and exporting stagflationary instability to the periphery \citep{Vernengo2021}, this hegemonic transition intersected with the secular post-war decline of corporate profitability across the advanced capitalist core \citep{Weisskopf1979,BowlesGordonWeisskopf1983,Vidal2019}. Yet the subsequent macroeconomic crisis in Chile cannot be reduced to a mechanical reflection of global monetary disorder. Rather, it represented the structural collision of ambitious domestic redistributive commitments with a targeted imperial credit squeeze, within which the Central Bank of Chile exercised constrained institutional agency to defend enterprise liquidity.

Mainstream and neoclassical syntheses characterize the Allende administration through the paradigm of \textit{``macroeconomic populism''} \citep{DornbuschEdwards1990,LarrainMeller1991}. In this interpretation, redistributive governments inevitably trace an unforced four-phase cycle: initial reactivation through idle capacity, the emergence of bottleneck shortages, unsustainable fiscal and monetary financing, and final economic collapse. By treating the central bank as an unconstrained printing press and monetary emission as an unmediated political choice, this literature suffers from two major analytical limitations. First, it succumbs to what \citet{BlancoMatute2018} formulate as the \emph{illusion of causality}---a cognitive heuristic that infers direct causal generation from the mere temporal coincidence of salient events (the 1970 election, 1971 wage readjustments, and 1972--1973 inflation) while ignoring non-contingent external covariates. Second, it exercises what \citet{IschEtAl2026} classify as \emph{narrative license}---the rhetorical escalation of observational and correlational snapshots into ironclad causal laws to satisfy academic demands for compelling monocausal parables. By analyzing aggregates in an institutional vacuum, mainstream accounts conflate correlation with causality, overlooking the methodological underpinnings of the ``credibility revolution'' and causal identification. Concretely, this conflation manifests itself by misidentifying the symptoms of redistribution, the targeted external financial blockade \citep{DeVylder1976,PetrasMorley1975}, and the balance-of-payments consequences of the contingent U.S. decision to close the convertibility window with the mechanical results of domestic monetary policy.

To discriminate between these contested claims on rigorous empirical ground, I deploy an applied Lakatosian evaluation of scientific research programmes in economic thought \citep{Lakatos1978,KingSznajder2006,Kvangraven2021}. This methodological framework evaluates rival programmes by distinguishing their foundational theoretical axioms (\emph{hard core}) from the auxiliary hypotheses of their empirical models (\emph{protective belt}), assessing whether empirical anomalies are accommodated through progressive theoretical extensions that predict \emph{novel facts} or through ad hoc adjustments. The debate over the macroeconomic destabilization of the Unidad Popular can be structured into three competing research programmes. The monetarist programme \citep{Edwards2024} anchors its hard core in exogenous base money creation, fractional-reserve multiplier causality, and long-run monetary neutrality in a closed system. Confronted with the empirical anomaly of 100\% monetary growth coinciding with non-inflationary output expansion in 1971, its protective belt relies on narrative rationalizations that dismiss the \$294~million external financing cutoff, the U.S. Eximbank embargo, and domestic distribution bottlenecks as secondary disturbances. The Austrian business-cycle programme \citep{EspinosaVejar2025} similarly emphasizes domestic distortions, positing that suppressed interest rates and state credit expansion generated intertemporal capital malinvestment and coordination breakdown.

In contrast, treating dependency theory as a research programme \citep{Kvangraven2021} generates excess empirical content by formalizing an open-economy balance-sheet mechanism within the programme's protective belt. Anchored in the structural realities of the international currency hierarchy and unequal access to world money \citep{Vasudevan2009,Basu2022}, this framework derives the external solvency constraint governing peripheral central banking. Using the Solvency Growth Gap ($\Delta s_{t-1} \equiv g^F_{t-1} - g^H_{t-1}$), the model identifies state-dependent regime shifts: when the central bank maintains adequate foreign exchange reserves, domestic liquidity expansion mobilizes idle capacity without igniting immediate inflation, and external commodity price shocks are buffered by the reserve buffer. However, once external credit embargoes and terms-of-trade deterioration push the monetary authority across the critical threshold into conditions of central bank insolvency ($\Delta s_{t-1} \le -4.615\%$), this insulation mechanism dissolves. The empirical evaluation reveals that consumer price shocks systematically precede base money creation, and reverse defensive accommodation exceeds forward monetary transmission in duration and cumulative response, compelling the central bank into liquidity injection to sustain enterprise payrolls under acute balance-of-payments distress.

This chapter deepens the theoretical foundations established in Chapter~2. Whereas Chapter~2 identified the external boundary of accumulation through an annual balance-of-payments threshold that throttles mechanization under unbalanced growth, this chapter uncovers the dynamic balance-sheet mechanisms through which that external constraint operates at high frequency. To overcome the temporal aggregation bias and closed-economy endogeneity pervasive in annual accounts \citep{DornbuschEdwards1990,Edwards2024}, my research design assembles a two-tier dataset at annual and monthly frequency time series from multiple secondary sources.  

The annual time-series dataset ($1920$--$2010$, $T=91$) draws functional income shares from \citet{Astorga2023}, balance-of-payments accounts from \citet{NievasPiketty2025}, and disaggregated net capital stocks from the Perpetual Inventory (GPIM) dataset developed in Chapter~2. At this annual frequency, the Austrian malinvestment hypothesis ($H_2$) fails to find empirical support: gross fixed investment in machinery and equipment contracted ($-2.3\%$) during the 1971 expansion while structural capacity utilization surged to an all-time peak of $\mu_{1971} = 98.44\%$, presenting empirical evidence consistent with reactivation rather than intertemporal capital misallocation. 

To isolate the transmission channels of the subsequent breakdown, the high-frequency monthly tier ($1960$--$1980$, $N=250$ continuous months) integrates primary Central Bank of Chile archival data accessed through the Application Programming Interface (API) of the \textit{Base de Datos Estadísticos} (BDE) of the Chilean Central Bank (BCCh), including Central Bank base monetary emission $H$ and $M1$, indices of gold, copper, and wholesale U.S. prices, and activity indicators gathered by \citet{DiazBahamonde2023} to build import capacity and structural imbalance indices. Historical international reserve statistics from the IMF provide monthly reserve assets not readily accessible elsewhere. Following the August 1971 closure of the U.S. gold convertibility window, my proxy of Prebisch's purchasing power capacity to import ($M^{\text{cap}}_t$) fell to $27.7$, driving the Real Structural Imbalance Ratio ($\Theta_t \equiv M_t / M^{\text{cap}}_t \times 100$) to an unsustainable peak of $283.4$ in May 1973. In the vector autoregressive and Threshold VAR systems, this physical import strangulation governs the regime switch: once foreign credit embargoes \citep{DeVylder1976,PetrasMorley1975} severely erode reserve backing and push the economy across the estimated Solvency Growth Gap threshold ($\Delta s_{t-1} \le -4.615\%$, entering the adverse solvency-growth state reflecting conditions of central bank insolvency), external price shocks penetrate domestic retail prices immediately, compelling defensive monetary accommodation to maintain enterprise payrolls amidst supply paralysis.

The remainder of this chapter is structured as follows. Section~\ref{sec:literature_review} surveys the contested traditions in Chilean political economy and formalizes the epistemological critique of orthodox causality. Section~\ref{sec:macro_framework} develops the theoretical framework of world money, financial subordination, and balance-sheet mechanics, formalizing how peripheral financial fragility and central-bank balance-sheet exposure condition macroeconomic transmission. Section~\ref{sec:data_sources} details the multi-frequency data compilation, archival sources, and variable construction, including the Central Bank solvency metric and structural import imbalance ratios. Section~\ref{sec:historical_empirical_results} establishes the historical background of the pre-1970 accumulation regime, presents stylized facts of peripheral accumulation and reserve depletion, and implements the econometric evaluation across the stationary vector-autoregressive Granger directional-precedence battery and the non-linear Threshold VAR. Section~\ref{sec:discussion_conclusion} discusses competing research programmes in the political economy of the Unidad Popular, evaluates the empirical findings alongside hegemonic contingency and central-bank agency, assesses causal identification limits, and outlines strategic policy implications. Five methodological and empirical appendices provide the balance-of-payments accounting derivation (Appendix~\ref{app:bop_levr}), the archival data codebook (Appendix~\ref{app:archival_codebook}), the multivariate PCMCI causal discovery results (Appendix~\ref{app:pcmci_ee1}), the complete 25-pathway Threshold VAR impulse response functions (Appendix~\ref{app:tvar_girf_atlas}), and the linear VAR diagnostic space (Appendix~\ref{app:linear_var_diagnostics}).

